\documentclass{article}

\usepackage[activate={true,nocompatibility},final,tracking=true,kerning=true,spacing=true,factor=1100,stretch=10,shrink=10]{microtype}
% activate={true,nocompatibility} - activate protrusion and expansion
% final - enable microtype; use "draft" to disable
% tracking=true, kerning=true, spacing=true - activate these techniques
% factor=1100 - add 10% to the protrusion amount (default is 1000)
% stretch=10, shrink=10 - reduce stretchability/shrinkability (default is 20/20)

\usepackage[utf8]{inputenc}
\usepackage{amssymb}
\usepackage{mathtools}
\usepackage{svg}
\usepackage{graphicx}
\usepackage[superscript, biblabel, nomove]{cite}
\usepackage{hyperref}
%\usepackage{todonotes}


%\usepackage{draftwatermark}
%\SetWatermarkText{DRAFT}
%\SetWatermarkScale{3}

\begin{document}

% Macros
\newcommand{\HAV}{\textsc{hav}}
\newcommand{\NOM}{\textsc{nom}}


\title{Havven: a stablecoin system\\ v0.3}
\author{Samuel Brooks, Anton Jurisevic, Kain Warwick}
\date{October 2017}


\begin{figure}
    \centering
    %\includesvg[width=0.33\textwidth]{img/block8logo}
    \includegraphics[width=0.33\textwidth]{img/havvenlogo}
\end{figure}
\maketitle

\hfill

\begin{abstract}
\noindent There is currently no effective decentralised unit of account. Previous attempts to create
stable tokens have either relied on significant centralisation or have been undermined by their
complexity. We present Havven, a representative money system which seeks to achieve price
stability with respect to an external asset. Havven is a dual-token solution, composed of a
stabilised exchange token and the reserve token which backs it. Users are incentivised to maintain
this distributed reserve, and to manage the stable token supply so that it is in proportion with the
value of the collateral. Because the collateral is encapsulated entirely within the system and
distributed among its users, we remove the need for a trusted central authority.
Such a stable cryptocurrency, useful for everyday economic purposes, will accelerate the adoption of
distributed ledger technology.
\end{abstract}

%\todo[inline]{Incorporate white paper criticisms.}

%\listoftodos{}

\pagebreak

\section{Introduction}

\subsection{Money and Cryptocurrencies}

\noindent The technology of money has three key functions: to act as a unit of account, a medium of
exchange and a store of value. In addition, money should ideally exhibit durability,
portability, divisibility, uniformity, limited supply, and acceptability.
As payment technology has advanced in recent years, it has become increasingly invisible
and it is often lost upon users of money that, like any technology, it can be improved.
Specifically, this means improving the performance of those desirable properties. \\

\noindent Bitcoin is an impressive technological advancement on existing forms of money because it simultaneously improves durability, portability, and divisibility. Further, it does so without requiring centralised control or the enforcement of a nation state from which to derive its value. This fixed monetary policy has protected Bitcoin from debasement and devaluation, allowing it to outperform other forms of money as a store of value. But this has created the potential for short-run volatility as Bitcoin lacks a mechanism to dynamically adjust to changing demand. \\

\noindent Bitcoin has thus tended to be a poor medium of exchange and an even worse unit of account.
In order for something to perform these functions it must remain
relatively stable against the price of goods and services over the short to medium term.

\subsection{Stablecoins}

\noindent A stablecoin is a cryptocurrency designed for price stability, such that it can function both as a medium of exchange and unit of account. It should ideally be as effective for making payments
as fiat currencies like the US Dollar, but still retain the desirable characteristics of Bitcoin, namely
transaction immutability, censorship resistance and decentralisation. \\

\noindent Cryptocurrencies that exhibit these characteristics are clearly a far better form of money; but adoption has been hindered by the volatility inherent in the inflexible monetary policies employed by decentralised systems. Thus stability continues to be one of the most valuable yet elusive characteristics for the technology. Clearly, efforts to create dynamic monetary policies within cryptoeconomic systems is still nascent, and significant research into stable monetary frameworks for cryptocurrencies is required. \\

\noindent The interested reader can also find additional discussion of stablecoins,
cryptoeconomics, competitors, and other related topics on our blog at \href{http://blog.havven.io}{\texttt{http://blog.havven.io}}.

\subsection{Havven}

\noindent Havven is a decentralised stability marketplace, where users requiring stability transact directly with those who provide the collateral to create it. This enables a novel form of representative money in which there is no requirement for a physical asset, thus removing problems of trust and custodianship. The asset used to back the stablecoin is a pool of reserve tokens that collectively represent the system itself; controlling these reserve tokens reflects participation in the Havven system, and they are a proxy for its value. Havven generates fees from users who transact in the stablecoin and distributes them among the holders of the reserve token, compensating them for underpinning the system. Havven therefore rewards those who actively participate in maintaining the stability of the system and charges those who benefit from its utility. These rewards are proportionally applied in response to the active management of the supply of the exchange token such that its price mirrors that of the asset it tracks. \\

\noindent Because we have created a system that generates cash flow for participants, we now have an asset which can be used as the collateral to support the stablecoin with a well-defined market value. The key to this is that the value of the system is measured in US dollars. This allows the system to issue a stablecoin which can be presented and redeemed for a percentage of the collateral tokens valued at \$1. Backing a stablecoin in this way is beneficial because such a cryptoeconomic system does not require trust in a centralised party; each participant has full transparency over how many tokens have been issued against the available collateral at all times. \\

\noindent The two linked tokens and the complex of incentives are described below: \\

\noindent \textbf{Havven}

\vspace{1mm}

\noindent The collateral token, whose supply is static.
The capitalisation of the havvens in the market reflects both the system's aggregate value and the reserve
which backs the stablecoin. Thus, users who hold havvens take on the role of maintaining stability.
Following Bitcoin, the Havven system will appear in upper case and singular; while the havven
token will be lower case and may be plural.

\vspace{2mm}

\noindent \textbf{Nomin}

\vspace{1mm}

\noindent The exchange token - the stablecoin - whose supply floats.
Its price as measured in fiat currency should be relatively stable.
Other than price stability, the system should also encourage some adequate level
of liquidity for nomins to act as a useful medium of exchange. \\

\noindent Each holder of havvens is able to issue a value of nomins in proportion to the dollar value
of the havvens they hold and are willing to place into escrow. If the user wishes to release their escrowed havvens, they must present the system with nomins in order to free their havvens and trade them again.
The holders of this token provide both collateral and liquidity, and in so doing assume some
level of risk. To compensate this risk, such nomin-issuers will be rewarded with fees the system levies
automatically as part of its normal operation. \\

\subsection{Design Rationale}

\noindent This issuance mechanism allows nomins to act as a form of representative money, where each nomin represents a share in the havven value held in reserve. Nomins derive value insofar as they provide a superior medium of exchange, and are effectively redeemable for a constant value of the denominating asset. In this paper, we use US dollars as this asset, but this could be any external and appropriately fungible asset, such as a commodity or a fiat currency. The system incentivises the issuance and destruction of nomins in response to changes in demand, to maintain a constant price. \\

%\noindent The design choice to back the system with a self-referential token was obvious;
%an asset-backed stablecoin with a cryptocurrency basket as reserve will always be inherently
%volatile, despite diversification, and will never be able to achieve the bespoke functionality
%of an asset which derives its value from stability. \\

\noindent In order to issue a nomin, there must be a greater value of havvens escrowed in the system. Overcollateralisation in this way provides confidence that nomins can be redeemed for their face value, even if there is a reduction in the value of the locked havvens. For example, if the collateralisation ratio is set to 100:1 then the price of havvens would need to fall 99\% for the nomin to no longer be overcollateralised. On the other hand, the ratio could be set to 1:1, but any decrease in havven price would mean the currency was no longer overcollateralised. In this situation, it would not be possible for all nomins to be redeemed for their face value, which may cause confidence in the system to drop.  \\

\noindent While overcollateralisation strongly supports a stable nomin price, it simultaneously reduces the efficiency of the system by limiting the utility of the nomin. A ratio of 1:1 provides far greater utility than 100:1, with respect to both the effiicency of the collateral and the revenue generated from transaction fees. There is clearly a trade-off between efficiency and resilience in Havven, which means the collateralisation incentives must therefore not only promote nomin price stability. They must also allow the system to operate as efficiently as possible, whilst maintaining resilience to reduction in the value of the collateral. The design of these incentives is described in detail in section 2. \\

\noindent Due to denominating the value of the stablecoin in an external fiat currency stability is relative only to that currency. In the future the system may support additional flavours of stablecoin that are denominated in other currencies such as the Euro. Denominating the stablecoin in a fiat currency removes the need to respond to macroeconomic conditions, as it benefits from the stabilisation efforts of large institutions acting in fiat markets. \\

\pagebreak

\section{Functional description}

Havven works by providing a set of market incentives that support the stability of nomin value with respect to an external asset.

\todo{Address how to maintain baseline demand.}
\todo{Address the role of the Havven foundation.}
\todo{Outline basic verbs for market players.}

\todo{Move analytical subsections from functional description section to the qualitative analysis and/or rationale sections.}
\subsection{Investment incentives}

We consider the reasons why any rational actor would buy curits. A potential buyer has at least three avenues for making money in Havven:

\paragraph{Capital gains due to the appreciation of curits:}
Due to its constrained supply, and the intrinsic utility of the stablecoin that it backs, it's reasonable to assume that
curits will appreciate in price.

\paragraph{Interest accrued from fees:}
If the price of curits stabilises for long periods of time, fees may be the only source of revenue. Ideally fees are set at a level where they are both high enough to be an incentive for rent-seekers to hold curits in the long term (thus assuming the risk of providing collateral for the system) and low enough not to be a disincentive for ordinary users to transact in nomins.
It is desirable, perhaps in a future world dominated by micropayments, for these fees to be negligible for end users, while still being macroeconomically important for the system, and for those who capitalise it.

\paragraph{Arbitrage profit:}
It is the arbitrageurs who will ultimately bring the price of nomins back into balance by a triangular circuit through nomins, curits, and the external (crypto or fiat) markets. Arbitrageurs might hold curits for a short time in order to pursue this strategy.


\pagebreak
\subsection{Fees}

There are several key considerations with respect to fee design:

\subsubsection{Fee design considerations}

\paragraph{The purpose of fees}

Fees are intended to be redistributed to actors who support the stability of the system. A fee pool will be distributed periodically for this purpose.
If the system determines that the nomin price is too low, then fees could be burned. If the price is too high then the system could sell these back
into the system at a discounted rate. The fee collection rate will also be a direct measure of the velocity of money in Havven. It's in the interest
of curit holders to maximise liquidity in order to maximise their return.

\paragraph{Fee beneficiaries}

One possibility would be simply to award fees to any holder of curits,
but in this situation holders can get all the benefit without taking any risk.
Although in the aggregate, it would be better for curit-holders if everyone issued nomins,
the marginal return for any single player (who cannot issue a large fraction of all circulating nomins)
of actually issuing them would not outweigh the risk they take on in doing so. If a user can issue 1\% of circulating
nomins, then doing so will only increase their fee takings by 1\%. Hence rational actors would not be incentivised to issue nomins at all.
This is a classic tragedy of the commons.

\noindent In order to avoid this situation, we must improve the marginal benefit of issuing nomins into circulation.
Hence, fees must be paid to those who \textit{issue} nomins, not just those who hold curits.

\paragraph{Fee collection}

The system can charge fees whenever any value is transferred, or any state is updated.
\noindent Different fee rates have different macroeconomic effects. We might in general like to set higher curit than nomin transfer fees, making the stablecoin itself a lower friction market in order to incentivise its use for exchange. Meanwhile, issuance and redemption fees will change the difficulty of entering and exiting the issuance game. \\

\noindent It is also possible for fees to float. The fee schedule could be altered dynamically in order to stabilise the system. It is even conceivable that the system could set negative fee rates if it needed to and charge punitive fees if a user is above the targeted utilisation ratio. For example, if nomin liquidity is low, meaning the system wants to incentivise issuance, then nomin transfer fees could increase, thus having the combined effect of increasing the interest accrued by issuers (thus incentivising issuance) and at the same time making it more expensive to transact in nomins. This would reduce demand and decrease the liquidity requirements. \\

\noindent Of note, fees are antithetical to arbitrage. The higher the fee, the higher the transaction friction, and the harder it is to make money by arbitrage. For example, if exchange fees amount to 1\% per trade, then a full arbitrage cycle between all three markets, (nomins, curits, and fiat) will cost in excess of 3\%. So it would not make sense to undertake arbitrage until such a time as the quoted exchange rate is misvalued by more than 3\% relative to the cross exchange rate. Hence, fees compete with arbitrage to stabilise price. Lower fees allow tighter stabilisation, within a window exactly in proportion with the fee rates themselves.

\subsection{Encouraging liquidity}

\noindent It's desirable that when actors issue nomins they are actually injected into the liquidity pool for their intended use,
rather than being held by the same actor in order to benefit from both the receipt of fees while retaining the option of using those nomins to rapidly release their curits.
In this manner they would accrue fees, but take on none of the risk of spending those nomins, for they always have an instantaneous option to liquidate their position and escape.
On the other hand, an actor who had done the economically-desirable thing and issued nomins to the market would be forced to buy them back before redeeming their escrowed curits.

\subsubsection{Non-discretionary Issuance}

One possibility is to simply provide an issuer no control over the tokens they issue. That is, when a quantity of nomins is issued, they are generated by the system which then places a sell order at the current going rate for that quantity on an exchange on the behalf of the issuer. When the order is filled, the proceeds in ether are remitted to the issuer. \\

\noindent Conversely, when a quantity of nomins is burned, they must first be obtained from the open market. In this way, a user would indicate an intention to burn, providing sufficient value to buy the proposed quantity of nomins, and the system would bid for that quantity on their behalf, thereby liquidating the user's curit position once the nomins have been obtained. \\

\noindent So one might consider there to be a formal distinction between wallets that issue tokens and those that do not. In this vein, one might envisage an extra fee to be charged to directly transfer nomins (rather than buying from the market) into a wallet that has an outstanding quantity of nomins it has previously issued, but not burnt. The result of this is that it would be less reasonable for an agent to sit on nomins in order to burn them in future as it is more advantageous in times of relative stability to simply buy them from the market. \\

\pagebreak 

\subsection{Price discovery}

One of the key challenges with denominating a cryptocurrency in a fiat currency is the
fundamental link this creates to the centralised world. When the denominating currency exists
external to the blockchain ecosystem, some bridge must be built so that the system can act with
knowledge of external information.
We recognise that a decentralised price-finding mechanism would be our preferred approach.
Research into this mechanism is on our horizon, and future versions of this white paper will contain
our results.
However, in order to reclaim system performance, we can trade some of the trustlessness of the design, for example by
a trusted ``Oracle'' service, which transmits knowledge of the external world into the system, building a causal link.


\subsection{Utilisation Ratio} Even though rational actor modelling suggests that the price of nomins and curits will equilibrilate
given that an agent may pay up to some multiple of the market value of a nomin in order to release escrowed curits, we are aware that there may be some
prevailing macroeconomic or psychological influences relating to an undercollateralised position
(i.e. if the value of the collateral pool is less than the issued stablecoin). As such, our modelling incorporates the notion of a ``utilisation ratio''
\(0 < U < 1\), such that the system is over-collateralised in an attempt to counteract these potential issues. It may be that resolving an
optimised utilisation ratio is beyond the ability of our agent-based modelling to determine, and as such, selecting this may need to be
informed by the activity of a live system. Thus it is currently intended for Havven to initially include in its governance model the power
to correct the utilisation ratio. This power can be removed over time as the system is proven, perhaps directly linked to some parametric milestones such as nomin velocity and stability.

%\paragraph{Direct Redemption} Direct redemption is a system design option to allow a holder of nomins to redeem any escrowed curits. This option allows more efficient redemption of nomins for the backing value, however introduces the difficulty of liquidating another actor's escrowed curits and interrupting their collection of fees. This can be solved by adding a premium to the price of an escrowed curit (potentially user-defined) over the current value. This premium would need to be paid in ether as to preserve the symmetric issuance and destruction of nomins. \\

\pagebreak

\section{Quantitative System Analysis}

We take the view that falsification is an important aspect of validating
the Havven system. In our quantitative analysis we seek to identify failure modes
of the system, and also to characterise not just \textit{whether} Havven stabilises
nomin prices, but \textit{how much} it does.

%\paragraph{Game-theoretic modelling}
% Discuss the structure of the game Havven represents, incorporating information the economists produce.
%\todo[inline]{Discuss game theoretic modelling.}

\paragraph{Agent-based modelling} It has been observed that analytic methods are often difficult to
apply in the complex and dynamic setting of a market.
One suggested solution to this problem is \textit{agent-based modelling}.
Under this paradigm, we proceed by first defining rational agent behaviour
and then simulating the interplay of those strategies over time.
We seek to develop a more effective
method of characterising market behaviour and equilibrium prices than pure analytic reasoning.~\cite{poggio2001agent}\\

\noindent Such simulations also provide an immediate means of measuring
quantities of interest. Simply by observing
the model, we can discover how varying input parameters
affect system outputs in an experimental fashion.
One important corollary is that this is a way of extracting reasonable
settings for system parameters (such as fee levels) that might be difficult
to reason about \textit{a priori}. These systems, reactive as they are,
also provide a method for testing proposed remedies for any identified failure
modes.

%\todo[inline]{Discuss that the game theory conclusions can be simulated}

\noindent In sum, then, the modelling seeks to answer the following, among other questions:

\begin{itemize}
    \item Does the system stabilise its nomin price?
    \item Under what conditions can stability fail?
    \item What are reasonable initial settings for fees and other parameters?
    \item What effect does the utilisation ratio have on havven/nomin price ratio?
    \item What is an effective utilisation ratio?
    \item What is the effect of a direct redemption regime?
    \item What are the expected returns for havven-holders?
\end{itemize}

%\todo[inline]{Fuller description of the technicals of the modelling.}

% \subsection{Modelling Havven}

% \paragraph{Environment}

% \paragraph{Agents}

% \paragraph{Optimal strategy mix}

\noindent Please visit \href{http://research.havven.io}{\texttt{research.havven.io}} for a pre-alpha version of our model.

\pagebreak

%
\section{Qualitative System Analysis}

This section provides a qualitative treatment of the reaction of the Havven system in response to various scenarios listed below:

\todo[inline]{Scenario analysis.}

\subsection{Incentives}

\paragraph{Why would anyone use nomins?}
\begin{itemize}
	\item Versus curits?
	\item Versus alternatives?
\end{itemize}

\paragraph{Why would anyone issue nomins?}
\begin{itemize}
	\item Fees
	\item Because they thought the peg would break in the positive direction.
\end{itemize}

\subsection{Scenarios}

\subsubsection{Curits appreciate against nomins.} \label{sec:curit-nomin-appreciation}
\begin{itemize}
	\item{Curit-holders can issue more nomins.}

	This might be scary because we either want the nomin supply to fall, or nomin demand to increase.
	This is fine if the nomin price fell, because then nomin demand should increase.
	However, if it was simply that the curit price increased, then this might encourage
	oversupply of nomins as issuers compete for shares of the fee pool.

	\item{Cheaper exit}

	Anyone who has previously issued curits but who wants to exit can buy nomins for
	cheaper than they issued them at, so egress with a profit.
	Alternatively, if the curit price doubles, then only half of their stake is required
	to back the nomins they have issued. They can use these proceeds to completely liquidate
	their position.

	\item{Each nomin locks fewer curits.}
	
\end{itemize}

\subsubsection{Curits depreciate against nomins.} \label{sec:curit-nomin-depreciation}

\begin{itemize}
	\item{Curit-holders can issue fewer nomins.}
	
	Perhaps they may even be under-staked.

	\item{Each nomin locks more curits.}
	\item{A player can issue a quantity of nomins and sell them for curits.}
		  
	They would do this on the assumption that, once the nomin price decreases as a result of the increased
	supply, they will be able to buy back the same quantity of nomins to free up their
	curits more cheaply.
\end{itemize}

\subsubsection{Nomin/curit liquidity dries up}
\begin{itemize}
	\item Low nomin supply

	Nomins should appreciate against the curit. See \cref{sec:curit-nomin-depreciation}.
	Fee takings will decrease, hurting curit-holders in the long run,
	which should incentivise them to inject nomins into the ecosystem.

	\item Low curit supply

	If system backers are sitting on a pile of curits, then the curit should
	appreciate. See \cref{sec:curit-nomin-appreciation,sec:curit-appreciation}.
\end{itemize}


\subsubsection{Long-run curit price appreciation}\label{sec:curit-appreciation}
\todo[inline]{Long-run curit appreciation scenario.}

\subsubsection{Long-run curit price depreciation}\label{sec:curit-depreciation}
\todo[inline]{Long-run curit depreciation scenario.}


\subsubsection{Radical shifts in usage}

\subsection{Expected Market Players}
\todo[inline]{List expected players in the market.}
\todo[inline]{Outline incentives and actions for different players.}

\begin{itemize}
	\item{Curit Holders}

	A curit-holder provides collateral and liquidity. It’s assumed curit-holders seek fee revenues,
	escrowing as many curits as they can. This incentive only really makes sense if curits are not
	significantly volatile over the long run. But in an unstable regime we also provide incentives
	for stabilising the nomin price. 

	\item{Nomin Users}

	Merchants, consumers, service providers, and so on: people who use the stable coin as a
	medium of exchange. They provide a base demand for nomins, which is necessary for fees
	to exist. These users may be disincentivised from using the system by excessive volatility
	in the price of nomins, or by high fees.

	\item{Arbitrageurs, Market Makers}

	The arbitrage force allows us to assume that the cur/nom, cur/fiat, nom/fiat
	prices are properly in alignment or will soon become aligned. Market making activities
	allow us to neglect the bid/ask spread, and situations where there is insufficient
	liquidity for players to transact.
	
	\item{Speculators}
	
	May tend to magnify price corrections, and are a significant vector by which to introduce
	exogenous shocks to the system modelling, e.g. large capital flows in response to breaking news.
	
	\item{Malicious Attackers}
	
	We should examine what happens if a George Soros (or otherwise) attacks Havven.
	\todo[inline]{List the various possible attacks against the system.}

	\item{Central Banker}
	
	The Havven foundation will have significant capital reserves with which it could intervene
	in the market if necessary to stabilise nomin prices. The system should work without
	such actions, but in extreme situations it might be necessary to undertake them.
	The advantage of such a market participant is, given that a very large market entity is
	willing to underwrite the stability of the coin, profit strategies predicated upon the
	stability of the token become less risky, so more feasible. So the Havven foundation in
	this capacity takes on the role of providing confidence.

\end{itemize}



\pagebreak

\section{System variables}

%\todo[inline]{More complete system variable section.}

\noindent What follows are the main variables of the system. Under each heading, each row will correspond to a single quantity of interest. Each row will have three columns. Leftmost, a mathematical definition of the variable; in the middle, the dimension of the quantity (which units it is measured in); and on the rightmost, a short English summary of the variable.\\

\noindent Certain abbreviations will be used. For example, \(\HAV{}\) and \(\NOM{}\) will be used as abbreviations for havvens and nomins considered as units of measurement. \\

\paragraph{Prices}
\begin{align*}
    P_v & \ && &(\frac{\text{\$}}{\HAV{}}) && &\text{: havven price.} \\
    P_n & \ && &(\frac{\text{\$}}{\NOM{}}) && &\text{: nomin price.} \\
    \pi &:= \frac{P_v}{P_n} \ && &(\frac{\NOM{}}{\HAV{}}) && &\text{: havven to nomin conversion factor.} \\
    P_v' &= f(V_n, V_v) \cdot R && &(\frac{\text{\$}}{\NOM{} \cdot \text{sec}}) && &\text{: havven price rate of change.}
    \intertext{Here \(R\) is a risk term incorporating, for example, volatility, number of buyers versus sellers, and so on.}
\end{align*}
\\


\paragraph{Money Supply}
\begin{align*}
    &H \ && &(\HAV{}) && &\text{: Quantity of havvens, which is constant.} \\
    &H_e \ && &(\HAV{}) && &\text{: Quantity of escrowed havvens.} \\
    &N = H_N \cdot \pi \ && &(\NOM{}) && &\text{: Quantity of nomins. This can float.} \\
    &H_N = \frac{N}{\pi} \ && &(\HAV{}) && &\text{: Havven value of issued nomins.}
    \intertext{Ideally, \(H_N \leq H_e\).}
\end{align*}
\\

\paragraph{Utilisation Ratios}
\begin{align*}
    &U = \frac{H_N}{H} \ && &\text{(dimensionless)} && &\text{: Empirical issuance ratio. } \\
    &U_{max} \ && &\text{(dimensionless)} && &\text{: Targeted issuance ratio ceiling.}
    \intertext{\(0 \leq U \leq U_{max} \leq 1\)}
\end{align*}
\\

\paragraph{Microeconomic Variables} These should be defined as functions of \(P_n, \ P_v, \ \text{fees, etc.}\)
\begin{align*}
S_n \ && (\frac{1}{\text{sec}}) && &\text{: average nomin spend rate} \\
S_i \ && (\frac{1}{\text{sec}}) && &\text{: average issuance rate} \\
S_r \ && (\frac{1}{\text{sec}}) && &\text{: average redemption rate}
\end{align*}
\\

\paragraph{Money Movement}
\begin{align*}
    V_n &= S_n \cdot N \ && &(\frac{\NOM{}}{\text{sec}}) && &\text{: nomin transfer rate.} \\
    V_v &= V_i + V_r \ && &(\frac{\HAV{}}{\text{sec}}) && &\text{: nomin} \leftrightarrow \text{havven conversion rate.} \\
    V_i &= (H - H_N) \cdot S_i \ && &(\frac{\HAV{}}{\text{sec}}) && &\text{: nomin issuance rate.} \\
    V_r &= H_N \cdot S_r \ && &(\frac{\HAV{}}{\text{sec}}) && &\text{: havven redemption rate.} \\
    \intertext{\(V_i\) is assumed to grow as there are more free havvens in the system.}
    \intertext{\(V_r\), by contrast, is taken to grow proportionally with the number of escrowed Havvens.}
\end{align*}
\\

\paragraph{Fees}
\begin{align*}
\intertext{The following fees are ratios, for example 0.1\%, levied on each transaction.}
&F_{nx} & \ && &\text{(dimensionless)} && &\text{: nomin transfer fee} \\
&F_{vx} & \ && &\text{(dimensionless)} && &\text{: havven transfer fee} \\
&F_i & \ && &\text{(dimensionless)} && &\text{: nomin issuance fee} \\
&F_r & \ && &\text{(dimensionless)} && &\text{: havven redemption fee} \\
% \intertext{These quantities are the aggregated fees accrued by the system per unit time.}
% &Ag_{nx} &:= V_n \cdot F_{nx} \ && &(\frac{\NOM{}}{\text{sec}}) && &\text{: fees taken from nomin transfers.}
\end{align*}

\pagebreak

\section{Alternative approaches}

\todo[inline]{Establish a summary of arguments against each competitor}
\todo[inline]{Makerdao critique.}
\subsection{Basecoin}

\paragraph{Description of system}

\noindent Basecoin is described as operating similarly to Havven in that there is separation between a backing token and a transactional token, however Basecoin also separates out a specific ``bond'' token.
The peg to an arbitrary external asset is maintained by using an oracle service to discover the price on an external market, before regulating the supply of ``basecoins'' through actively increasing (issuing new basecoin),
and decreasing (auctioning of bonds) the supply, effectively acting as an autonomous central bank.

\subsubsection{Key issues}

\noindent Basecoin is intended to operate ``as a decentralized, protocol-enforced algorithm, without the need for direct human judgment (sic). For this reason,
Basecoin can be understood as implementing an algorithmic central bank.'' Whilst not without merit, this approach was discarded by Havven due to the high degree
of design complexity required to be anticipated in order to ensure the stabilisation mechanism is effective. The paper states that Monte Carlo simulations have
been run which indicate stability under a range of scenarios, however details are yet to be released by the team. \\

\noindent Another element not explored in the Basecoin whitepaper is the incentives for participants to engage with the cryptoeconomic system itself.
While there is no argument against the utility of stablecoins, there must be incentives inherent in all such systems to ensure the appropriate
participation of all actors. In this case, there are consumers of the stablecoin and active participants in the monetary policy. It is critical
to be able to demonstrate that the incentives within the system will ensure profitable participation strategies for actors. Without this being clarified
it is unclear as to whether there will be uptake by enough users to generate sufficient currency in circulation to support the demand for a stablecoin.
Critically, the removal of Basecoin from the system to ensure the stable peg is predicated on the significant assumption that participants will take positions
in the ongoing bond auctions. This assumption remains untested. \\

\noindent A final point needs to be made with respect to the overarching monetary approach
espoused in the whitepaper. In the section ``Averting Macroeconomic Depressions'' the authors appear to support
money printing and inflationary policies and the subsequent devaluation of currency. Even were it possible to
demonstrate that inflation of the money supply via such a system would be effective in combating a deflationary spiral,
a far better argument could be made that simply by implementing a stable store of value and unit of account that such a
system would not be required. Generally, the apparent assumption that such a system would be achievable and still able to
handle monetary crises in a far future time without centralised intervention stretches credulity. It's not entirely clear
why Basecoin has intended to merely replicate the function of a central bank, rather than aim for pure stability or a
relative-stable approach such as Havven. We are skeptical of any group that would advocate for monetary approaches that are
diametrically opposed to cryptoeconomic efforts to democratise money, and we feel that the proposal to intentionally
create a systematically inflationary monetary system is not the answer. Instead, we should at this point in time be
aiming to construct a system that provides a stable store of value relative to an arbitrary fiat currency. The macroeconomic
benefits of such a system are clear, and for as long as we live in a fiat-dominated world this will continue to be the case.

\subsection{Tether}

\paragraph{Description of system}

Tethers accepts fiat deposits into the Hong Kong-based Tether Limited bank account and issues ``USDT'' (USD Tether) over Bitcoin via the Omni Layer protocol. Tethers are an asset-backed digital token, representing a claim on the cash held in reserve. \\

\noindent The stability of the USDT `coin' effectively relies on the force of external market arbitrage to ensure the peg holds over time.

\paragraph{Key issues}

Despite the whitepaper claiming that the ``goal of any successful cryptocurrency is to completely eliminate the requirement for trust,'' and that each Tether is ``fully redeemable/exchangeable any time for the underlying fiat currency,'' the company's terms of service quite clearly state that ``there is no contractual right or other right or legal claim against us to redeem or exchange your Tethers for money.'' \\

\noindent Tether clearly relies on a manual, centralised proof of existence for the backing asset, and so suffers from the very issue that the Tether whitepaper decries. Indeed the same issue is encountered with tokenised gold, or similarly any other real-world asset where some Oracle bridge is required to interface into a distributed ledger.

\paragraph{Current state}

Recently, Tether announced support for issuing ERC-20 compatible tokens on Ethereum as opposed to releasing ``tethers'' on the Bitcoin blockchain using the Omni Layer protocol. \\

\noindent At the time of writing, the market capitalisation for USDT was approximately \$440m, and the discrepancy regarding their terms of service remains unresolved. \\


% \subsection{MakerDAO}

% \paragraph{Description of system}

% \paragraph{Key issues}

% \paragraph{Current state}

% \subsection{Nubits}

% \paragraph{Description of system}

% \paragraph{Key issues}

% \paragraph{Current state}

\pagebreak

\bibliography{tex/citations}
\bibliographystyle{plain}

\end{document}
